\subsection{Preliminaries}\label{rl-0003}
Group actions on algebraic stacks are subtle, and there are several distinct notions of ``the fixed stack'' for a group acting on a stack. We make no attempt at a comprhensive treatment here; instead, we review 

We begin by reviewing group actions on stacks as defined by Romagny \cite{romagny_GroupActionsStacks2005} and four distinct notions of ``the fixed stack of a group action on a stack''. Throughout the main body of this paper we use the Romagny fixed stack

\begin{definition}\label{rl-0004}
\cite[Definition 2.1]{romagny_GroupActionsStacks2005}
  \label{defn:group-action-on-a-stack}
  Let \(\mathcal X\) be a stack over a scheme \(S\) and let \(G\) be a group algebraic space. Let \(m:G\times G\to G\) be the multiplication of \(G\) and \(e:S\to G\) its unit section. Let \(U\) be an \(S\)-scheme.
  \begin{enumerate}[(i)]
    \item An \textit{action} of \(G\) on \(\mathcal X\) is a morphism of stacks \(\mu:G\times \mathcal X\to \mathcal X\) satisfying the following 1-commutative diagrams
          \begin{equation}\label{eqn:group-action-diagrams}
            \begin{tikzcd}
            G\times G\times \mathcal X \arrow[r, "m\times \id_{\mathcal X}"] \arrow[d, swap, "\id_G\times \mu"] & G\times\mathcal X \arrow[d, "\mu"] \\
            G\times \mathcal X \arrow[r, "\mu"] & \mathcal X
            \end{tikzcd}\hphantom{asdfasdf}
            \begin{tikzcd}
            G\times \mathcal X \arrow[r,"\mu"] & \mathcal X \\
            \mathcal X\arrow[u, "e\times \id_{\mathcal X}"] \arrow[ur, swap, "\id_{\mathcal X}"{name=D}] &
            \end{tikzcd}
          \end{equation}
          The pair \((\mathcal X, \mu)\) is called a \(G\)-stack.
    \item A \(1\)-morphism of \(G\)-stacks or a \emph{\(G\)-equivariant morphism} between \((\mathcal X, \mu)\) and \((\mathcal Y, \nu)\) is a pair \((f, \sigma)\) where \(f:\mathcal X\to \mathcal Y\) is a \(1\)-morphism and \(\sigma:\nu\circ (\id_G\times f)\to f\circ \mu\) is a \(2\)-isomorphism written on sections as \(\sigma^x_g:g\cdot f(x) \xrightarrow{\sim} f(g\cdot x)\), such that \(\sigma_g^{h\cdot x}\circ g \cdot \sigma^x_h = \sigma^x_{gh}\) for all \(U\), \(x\in \mathcal X(U)\) and \(g, h\in G(U)\).
    \item A \textit{2-morphism} of \(G\)-stacks between 1-morphisms \((f_1,\sigma_1)\) and \((f_2, \sigma_2)\) is a 2-morphism \(\tau:f_1\Rightarrow f_2\) so that for all \(x\in \mathcal X(U)\) and \(g\in G(U)\) we have \(\sigma_{2,(g,x)}\circ g\cdot \tau_x = \tau_{g\cdot x}\circ \sigma_{1,(g,x)}\); by \(g\cdot \tau_x\) here we mean \(\nu(g, \tau_x)\).
    \item An \textit{isomorphism} of \(G\)-stacks is a \(1\)-G-morphism that is also an equivalence of groupoids over \(S\). A monomorphism is a fully faithful \(1\)-G-morphism and an epimorphism is a \(1\)-G-morphism that is locally essentially surjective.
  \end{enumerate}
\end{definition}

\begin{remark}\label{rl-0005}
  Technically what is stated here is what Romagny calls a \emph{strict action}; the more general version asks that these diagrams only commute up to 2-isomorphism. However Romagny immediately shows that strictification is a 2-equivalence \cite[Proposition 1.5]{romagny_GroupActionsStacks2005}
\end{remark}


Consider for a moment the fixed locus \(X^G\) of a scheme \(X\) with a group action \(G\). A map \(U\to X^G \in X^G(U)\) from some scheme \(U\) is exactly the same data as a \(G\)-equivariant map \(U\to X\) where \(U\) is endowed with the trivial action of \(G\). Romagny defines the fixed stack analogously to this functor of points description of \(X^G\). 
\begin{definition}\label{rl-0006}
  \label{defn:romagny-fixed-stack}
  Let \(G\) be a group algebraic space over \(S\) a scheme and let \(\mathcal X\) be a \(G\)-stack over \(S\). For an \(S\)-stack \(Y\), denote by \(\iota(\mathcal Y)\) the \(G\)-stack over \(S\) obtained by endowing \(\mathcal Y\) with the trivial \(G\)-action. The \textit{stack of fixed points} \(\mathcal X^{G}\) is a stack that \(2\)-represents the functor
  \[\sHom_{G-\St}(\iota(-), \mathcal X).\]
  Note that \(\mathcal X^G\) is an object in the category of stacks over \(S\), and \textit{not} in the category of \(G\)-stacks over \(S\). There is a canonical morphism \(\eta:\mathcal X^G\to \mathcal X\) which sends \((f,\sigma)\in \mathcal X^G(U)\) to \(f\in \mathcal X(U)\), where here we consider \(\mathcal X\) without its \(G\)-stack structure. 
\end{definition}

\begin{remark}\label{rl-0007}\label{rmk:description-of-homotopy-fixed-stack}
  Here is a slightly more explicit description of the points of $\mathcal X^{G}$ over a scheme $U \in \Sch/S$. An object of $\mathcal X^{G}(U)$ consists of a $G$-equivariant morphism $(x,\alpha):\iota(U)\to X$. The equivariance diagram in this case is
  \begin{center}
    \begin{tikzcd}[column sep={1.5cm,between origins},
      row sep={1.5cm,between origins}]
      G\times U\arrow[r, "\pr_2"] \arrow[d, swap, "\id_G \times x"] & U \arrow[d, "x"] \\
      G\times \mathcal X\arrow[r, "\mu"] \arrow[ur, Rightarrow, "\alpha"]& \mathcal X
    \end{tikzcd}
  \end{center}
  since $\iota(U)$ is simply $U$ with the trivial $G$-action. The morphism $x$ is an object of \(\mathcal X\) by the Yoneda lemma and the 2-commutativity implies that for all $g \in G(U)$ we get an isomorphism $\alpha_g:g\cdot x\xrightarrow{\sim}x$. Similarly, given a pair $(x,\{\alpha_g\}_{g\in G(U)})$ consisting of an object $x \in \mathcal X(U)$ and a collection of isomorphisms $\alpha_g:gx\xrightarrow{\sim}x$, we obtain a $G$-equivariant map $\iota(U)\to \mathcal X$. Thus the objects of $\mathcal X^{G}(U)$ are precisely pairs $(x,\{\alpha_g\}_{g\in G(U)})$ where $x \in \mathcal X(U)$ and the $\alpha_g:g\cdot x\xrightarrow{\sim} x$ are isomorphisms satisfying the cocycle condition \(\alpha_{gh} = \alpha_g\circ g\cdot \alpha_h\) pictured by
  \begin{equation}\label{eqn:cocycle-condition-for-fixed-stack}
    \begin{tikzcd}[column sep={1.5cm,between origins},
      row sep={1.5cm,between origins}]
      & gx \arrow[dr,"\alpha_g"] & \\
      (gh)x\arrow[ur, "g\cdot \alpha_h"] \arrow[rr, "\alpha_{gh}"] & & x
    \end{tikzcd}
  \end{equation}
  for all $g, h\in G(U)$. The morphism \(\eta:\mathcal X^G\to \mathcal X\) simply forgets the data of \(\alpha\).

  In particular, when \(G\) acts trivially on \(\mathcal X\), \(gx = x\) for all \(x\in \mathcal X(U)\) and the cocycle condition reduces to the requirement that \(\alpha_{gh} = \alpha_g\circ \alpha_h\); i.e. that \(\alpha:G\to \underline{\Aut}(x)\) is a group homomorphism. Thus, for the trivial action, the fiber of \(x\in \mathcal X(U)\) over \(\eta\) is the sheaf \(\underline{\Hom}_{S-gp}(G, \Aut(x))\).
\end{remark}


When there are many choices of \(\alpha_g\) for a single \(x\in \mathcal X(U)\) we get multiple points in \(\mathcal X^G(U)\) mapping to a single point in \(\mathcal X\). This means that \(\eta:\mathcal X^G\to \mathcal X\) is, in general, not a monomorphism.

\begin{example}\label{rl-0008}
  Let \(\mathfrak M\) denote the stack of nodal curves of some fixed genus (we suppress little \(g\) from the notation since we want to use it to denote a group element) and endow it with the trivial action of an algebraic group \(G\). Consider a scheme \(B\). Specializing Remark \ref{rmk:description-of-homotopy-fixed-stack} to this situation, we see that the points of \(\mathfrak M^G(B)\) are given by pairs \((x,\alpha)\) where \(x\in \mathfrak M(B)\) and \(\alpha:G(B)\to \Aut(x)\) is a group homomorphism. The point \(x\) defines a family of nodal curves \(E\to B\) over \(B\) by pulling back the universal curve over \(\mathfrak M\), and \(\alpha\) gives an automorphism
  \begin{center}
    \begin{tikzcd}
      E\arrow[r, "\sim"] \arrow[rd] & E \arrow[d]\\
      & B
    \end{tikzcd}
  \end{center}
  of this family for every element of \(G(B)\). Thus, \(\mathfrak M^G\) parameterizes families of nodal curves with an action of \(G_B\) on \(E\) over \(B\), and \(\eta:\mathfrak M^G\to \mathfrak M\) is the map that forgets this \(G\)-action. Whenever there is such a family \(E\to B\) with a nontrivial \(G_B\)-action on \(E\), \(\eta\) will not be a monomorphism.
\end{example}


\begin{example}\label{rl-0009}\label{example:BGm-and-its-fixed-stack}
  Let \(S\) be a scheme with a trivial \(\mathbb G_m\) action, and endow \(\mathcal X = [S/\mathbb G_m] = B\mathbb G_m\) with the trivial \(T = (\mathbb G_m)^n\) action. Every point in \(x \in \mathcal X\) is fixed by \(T\), so a collection \(\{\alpha_t\}\) satisfying the cocycle condition \eqref{eqn:cocycle-condition-for-fixed-stack} is the same data as a homomorphism \(\alpha:T\to \Aut_{\mathcal X}(x) = \mathbb G_m\), i.e. a character of the torus. The fixed stack \(\mathcal X^T\) is therefore a countable union of copies of \(\mathcal X\):
  \[\mathcal X^T = \coprod_{\chi\in \mathbb Z^n} \mathcal X,\]
  and the fiber of \(\eta\) at a point \(x \in \mathcal X\) is \(\mathbb Z^n\). In particular, \(\eta\) is not a monomorphism and is not quasi-compact. The pushforward \(\eta_*\), the thing which would becomes an isomorphism in the classical localization, is not even defined.
\end{example}


The stack \(\mathcal X^G\) also admits a description via something called the \emph{universal stabilizer} of the action of $G$ on $\mathcal X$, denoted \(\St_{X,G}\) \cite[Section 4.1]{romagny_AlgebraicitySmoothnessFixed2022}. It is the fiber product
\begin{equation}\label{eqn:universal-stabilizer-diagram}
\begin{tikzcd}
  \St_{\mathcal X,G} \arrow[d]\arrow[r]& \mathcal X \arrow[d,"\Delta"] \\
  G\times_S \mathcal X\arrow[r,swap,"\text{act}\times \pr_2"]&\mathcal X\times \mathcal X
\end{tikzcd}.
\end{equation}
Over a point \(x:U\to \mathcal X\) its objects are
\begin{align*}
  \St_{\mathcal X,G}(U) = \{(g,\alpha) ~\mid~ g\in G(U), \alpha:gx\xrightarrow{\sim} x\},
\end{align*}
with multiplication law \((g,\alpha) \cdot (h,\beta) = (gh, \alpha \circ g\beta)\) and identity element \((1,\id_x)\). The leftmost vertical arrow in diagram \eqref{eqn:universal-stabilizer-diagram} is given by \((g,\alpha) \mapsto g\) over \(x:U\to \mathcal X\) and is a morphism of \(\mathcal X\)-group stacks. Its kernel is the inertia stack \(I_{\mathcal X} := \mathcal X\times_{\mathcal X\times \mathcal X}\mathcal X,\) and thus we get an exact sequence
\begin{equation}
  \label{eqn:inertia-stack-sequence}
  1 \to I_{\mathcal X/S} \to \St_{\mathcal X,G} \to G\times_S \mathcal X
\end{equation}
where over \(x:U\to \mathcal X\) the map \(I_{\mathcal X/S}\to \St_{\mathcal X, G}\) is
\begin{align*}
  (\alpha:x\xrightarrow{\sim} x) \mapsto (1, \alpha).
\end{align*}
We call the cokernel of \(I_{\mathcal X/S}(x)\hookrightarrow \St_{\mathcal X, G}(x)\) the \emph{\(G\)-stabilizer} of \(x\). As an aside, the rightmost map in \eqref{eqn:inertia-stack-sequence} admits a section over \(U\) if and only if \(x\) is in the image of \(\eta(U):\mathcal X^G(U)\to \mathcal X(U)\), and \(x:U\to \mathcal X\) in the image of \(\eta\).
\begin{remark}\label{rl-000A}
\end{remark}



\begin{lemma}\label{rl-000B}
\cite[Lemma 4.1.2]{romagny_AlgebraicitySmoothnessFixed2022}
  \label{lem:fixed-stack-parameterizes-splittings}
  $\mathcal X^{G}$ parameterizes sections of \(\St_{\mathcal X,G}\to G\times_S \mathcal X\).
\end{lemma}
\begin{proof}
  Fix a point \(x:U\to \mathcal X\). A splitting of \(\St_{\mathcal X, G}\times_{\mathcal X}U\to G\times_S\mathcal X\times_{\mathcal X}U\) is determined by a collection \(\{\alpha_g\}_{g\in G(U)}\) satisfying the cocycle condition in diagram \eqref{eqn:cocycle-condition-for-fixed-stack}, which is exactly the data of a point in \(\mathcal X^G(U)\) over \(x\).

  In particular, the image of \(\eta(U):\mathcal X^G(U)\to \mathcal X(U)\) is the collection of points \(x\in \mathcal X(U)\) where sequence \eqref{eqn:inertia-stack-sequence}, base changed to \(U\), admits a group-theoretic section over \(U\).
\end{proof}

The fixed stack \(\mathcal X^G\) and the map \(\eta:\mathcal X^G\to \mathcal X\) can be arbitrarily bad, but crucially, it is controllable for algebraic tori acting on algebraic stacks.
